\SourceLeaf{430}%
Fläche selbst gelegt wird. Für den Krümmungsradius haben wir also die einfache Formel
\[
\frac1r=Z(\xi\xi T+2\xi\eta U+\eta\eta V).
\]

\Needspace{5\baselineskip}
\begin{center}11.\end{center}
\noindent\hspace*{1.2em}Da sich durch diese Normale unendlich viele Ebenen legen lassen, kann der Krümmungsradius demnach unendlich viele verschiedene Werte annehmen. Hierbei sind $T,U,V,Z$ als konstant und $\xi,\eta,\zeta$ als veränderlich zu betrachten. Um Letztere von einer einzigen veränderlichen Größe abhängig zu machen, nehme man auf dem Großkreis, dessen Pol $L$ ist, zwei feste, um $90^{0}$ voneinander entfernte Punkte an, die $M,M'$ heißen mögen. Ihre Koordinaten in Bezug auf den Mittelpunkt der Kugel seien $\alpha,\beta,\gamma$; $\alpha',\beta',\gamma'$. Dann gilt
\[
\cos\lambda(1)=\cos\lambda M.\cos M(1)+\cos\lambda M'.\cos M'(1)+\cos\lambda L.\cos L(1).
\]
Setzt man nun
\[
\lambda M=\varphi,
\]
so gilt
\[
\cos\lambda M'=\sin\varphi
\]
und die Formel wird zu
\[
\xi=\alpha\cos\varphi+\alpha'\sin\varphi
\]
und ebenso
\[
\begin{aligned}
\eta&=\beta\cos\varphi+\beta'\sin\varphi\\
\zeta&=\gamma\cos\varphi+\gamma'\sin\varphi.
\end{aligned}
\]

Setzt man also
\[
\begin{aligned}
A&=(\alpha\alpha T+2\alpha\beta U+\beta\beta V)Z\\
B&=(\alpha\alpha'T+(\alpha\beta'+\alpha\beta')U+\beta\beta'V)Z\\
C&=(\alpha'\alpha'T+2\alpha'\beta'U+\beta'\beta'V)Z
\end{aligned}
\]
so gilt:
